\documentclass[11pt]{article}
\usepackage[margin=1in]{geometry}
\usepackage{amsmath,amssymb,amsthm}
\usepackage{algorithm}
\usepackage{algpseudocode}
\usepackage{booktabs}
\usepackage{hyperref}

\newtheorem{definition}{Definition}
\newtheorem{theorem}{Theorem}
\newtheorem{example}{Example}

\title{Epsilon-Approximate Credal Variable Elimination:\\
A Fully Polynomial-Time Approximation Scheme}
\author{IBM Research}
\date{}

\begin{document}
\maketitle

\begin{abstract}
The exact Credal Variable Elimination (CVE) algorithm may produce intermediate
potentials whose cardinality grows exponentially. By relaxing the pruning criterion
to allow \emph{epsilon-dominated} functions to be removed, we obtain an approximate
algorithm whose error is bounded by a user-specified tolerance $\varepsilon > 0$.
This epsilon-approximate CVE constitutes a fully polynomial-time approximation scheme
(FPTAS) for marginal inference in credal networks of bounded treewidth and bounded
variable cardinality.
\end{abstract}

\section{Motivation}

In exact CVE (see companion document \texttt{credal\_ve.tex}), the pruning step removes
only strictly dominated functions: $f$ is pruned if there exists $g$ such that
$g(\mathbf{y}) \geq f(\mathbf{y})$ for all $\mathbf{y}$, with strict inequality
somewhere. While this preserves exact bounds, the number of surviving functions can
grow exponentially across elimination steps.

The key insight of Mau\'a et al.\ (2012) is to relax the dominance criterion: remove
functions that are \emph{almost} dominated, controlled by a tolerance parameter
$\varepsilon$. This trades bound tightness for computational efficiency.

\section{Epsilon-Dominance}

\begin{definition}[$\varepsilon$-Dominance]
Let $\varepsilon \geq 0$. A function $f$ is \textbf{$\varepsilon$-dominated} by
function $g$ if:
\begin{equation}\label{eq:eps-dom}
  g(\mathbf{y}) \geq f(\mathbf{y}) - \varepsilon \quad \forall\, \mathbf{y}.
\end{equation}
\end{definition}

When $\varepsilon = 0$, this reduces to standard dominance. For $\varepsilon > 0$,
more functions are classified as dominated and thus pruned, reducing the potential's
cardinality.

\begin{definition}[$\varepsilon$-Pruning]
The $\varepsilon$-pruning of a potential $\phi$ is:
\[
  \mathrm{prune}_\varepsilon(\phi) = \{f \in \phi : \nexists\, g \in \phi \setminus \{f\}
  \text{ s.t.\ } g(\mathbf{y}) \geq f(\mathbf{y}) - \varepsilon\; \forall\, \mathbf{y}\}.
\]
\end{definition}

\section{Algorithm}

The algorithm is identical to exact CVE (Algorithm~\ref{alg:cve-eps}) except that
the pruning step at line~\ref{line:prune} uses $\varepsilon$-pruning instead of
exact pruning.

\begin{algorithm}[H]
\caption{Epsilon-Approximate Credal Variable Elimination}
\label{alg:cve-eps}
\begin{algorithmic}[1]
\Require Extreme points $\mathrm{ext}\,K(X_i \mid \mathrm{Pa}(X_i))$ for each node,
         query $Z$, evidence $\mathbf{Y} = \mathbf{y}$, tolerance $\varepsilon > 0$
\Ensure Approximate bounds $[\underline{P}_\varepsilon,\; \overline{P}_\varepsilon]$
\Statex
\State Initialize potentials $\Gamma$ from extreme points and evidence indicators
       \Comment{Same as exact CVE}
\State Compute elimination ordering $\mathbf{o}$ excluding $Z$
\For{each $X_i$ in ordering $\mathbf{o}$}
  \State $\Gamma_{X_i} \gets \{\phi \in \Gamma : X_i \in \mathrm{scope}(\phi)\}$
  \State $\Gamma \gets \Gamma \setminus \Gamma_{X_i}$
  \State $\psi \gets \prod \{\phi \in \Gamma_{X_i}\}$
  \State $\psi' \gets \sum_{X_i} \psi$
  \State $\psi'' \gets \mathrm{prune}_\varepsilon(\psi')$
  \label{line:prune}
  \Comment{$\varepsilon$-pruning instead of exact pruning}
  \State $\Gamma \gets \Gamma \cup \{\psi''\}$
\EndFor
\State Extract bounds from remaining potentials
\Comment{Same as exact CVE}
\State \Return $[\underline{P}_\varepsilon,\; \overline{P}_\varepsilon]$
\end{algorithmic}
\end{algorithm}

\subsection{Epsilon-Pruning Procedure}

\begin{algorithm}[H]
\caption{$\varepsilon$-Prune}
\label{alg:eps-prune}
\begin{algorithmic}[1]
\Require Potential $\phi = \{f_1, \dots, f_m\}$, tolerance $\varepsilon$
\Ensure Pruned potential $\phi'$
\State $\mathrm{dominated}[i] \gets \mathsf{false}$ for $i = 1, \dots, m$
\For{$i = 1, \dots, m$}
  \If{$\mathrm{dominated}[i]$} \textbf{continue} \EndIf
  \For{$j = 1, \dots, m$, $j \neq i$}
    \If{$\mathrm{dominated}[j]$} \textbf{continue} \EndIf
    \If{$f_j(\mathbf{y}) \geq f_i(\mathbf{y}) - \varepsilon$ for all $\mathbf{y}$}
      \State $\mathrm{dominated}[i] \gets \mathsf{true}$
      \State \textbf{break}
    \EndIf
  \EndFor
\EndFor
\State \Return $\phi' = \{f_i : \neg\,\mathrm{dominated}[i]\}$
\end{algorithmic}
\end{algorithm}

\section{Theoretical Properties}

\begin{theorem}[FPTAS, Mau\'a et al.\ 2012]
For credal networks of bounded treewidth $w^*$ and bounded variable cardinality $k$,
the $\varepsilon$-approximate CVE runs in time polynomial in the input size and in
$1/\varepsilon$, and produces bounds within $\varepsilon$ of the exact bounds:
\[
  \underline{P}(z \mid \mathbf{y}) \leq \underline{P}_\varepsilon(z \mid \mathbf{y})
  \leq \underline{P}(z \mid \mathbf{y}) + O(\varepsilon),
\]
and similarly for the upper bound.
\end{theorem}

\paragraph{Bound direction.}
Since $\varepsilon$-pruning only removes functions (never adds them), the approximate
algorithm produces bounds that are at least as wide as the exact bounds:
\[
  \underline{P}_\varepsilon \leq \underline{P}, \qquad
  \overline{P}_\varepsilon \geq \overline{P}.
\]
Thus the $\varepsilon$-approximate bounds are always valid \emph{outer} approximations.

\section{Running Example}

\begin{example}[Alarm Network]
We use the same alarm network as in \texttt{credal\_ve.tex} with variables $B$, $E$,
$A$, and $CD$.

\paragraph{Query: $P(B{=}1)$, elimination order: $E, A, CD$.}

\begin{center}
\begin{tabular}{lccc}
\toprule
$\varepsilon$ & $\underline{P}(B{=}1)$ & $\overline{P}(B{=}1)$ & Time (s) \\
\midrule
$0$ (exact) & 0.1000 & 0.2000 & 3.69 \\
$0.001$ & 0.1000 & 0.2000 & 3.14 \\
$0.01$ & 0.1000 & 0.2000 & 1.53 \\
$0.1$ & 0.1000 & 0.2000 & 0.009 \\
\bottomrule
\end{tabular}
\end{center}

For this query, even aggressive pruning ($\varepsilon = 0.1$) produces the same bounds
because $B$ is a root node with only two vertices. The speedup comes from pruning
functions in the intermediate potentials during elimination of $E$, $A$, and $CD$.

\paragraph{Query: $P(A{=}1 \mid B{=}0, E{=}0)$.}

\begin{center}
\begin{tabular}{lccc}
\toprule
$\varepsilon$ & $\underline{P}(A{=}1 \mid \cdot)$ & $\overline{P}(A{=}1 \mid \cdot)$ & Time (s) \\
\midrule
$0$ (exact) & 0.0000 & 0.0444 & 0.012 \\
$0.01$ & 0.0000 & 0.0444 & 0.002 \\
$0.1$ & 0.0000 & 0.0000 & 0.001 \\
\bottomrule
\end{tabular}
\end{center}

With $\varepsilon = 0.1$, the interval collapses to $[0, 0]$ because the true upper
bound ($0.044$) is smaller than $\varepsilon$, so the function achieving the maximum
gets $\varepsilon$-dominated. This illustrates the tradeoff: large $\varepsilon$
gives speed but may lose precision when probabilities are small relative to $\varepsilon$.
\end{example}

\section{Practical Guidance}

\begin{itemize}
  \item Start with $\varepsilon = 0$ (exact) for small networks.
  \item For larger networks, use $\varepsilon = 10^{-3}$ to $10^{-2}$ as a good
        balance between speed and accuracy.
  \item Avoid $\varepsilon$ larger than the smallest probability of interest.
  \item The algorithm is most effective when the network treewidth is moderate (the
        intermediate potentials are the bottleneck, and pruning directly controls
        their size).
\end{itemize}

\section{References}

\begin{enumerate}
\item D.D.\ Mau\'a, C.P.\ de Campos, M.\ Zaffalon. Solving limited memory influence
      diagrams. \emph{JAIR}, 44:97--140, 2012.
\item D.D.\ Mau\'a, F.G.\ Cozman. Thirty years of credal networks: Specification,
      algorithms and complexity. \emph{IJAR}, 126:133--157, 2020.
\item R.\ Marinescu et al. Credal Marginal MAP. \emph{NeurIPS}, 2023.
\end{enumerate}

\end{document}
